\documentclass[11pt]{article}
\usepackage[a4paper,margin=2.6cm]{geometry}
\usepackage{amsmath,amssymb,amsthm,mathtools}
\usepackage{microtype}
\usepackage[hidelinks]{hyperref}

\newtheorem{theorem}{Theorem}
\newtheorem{lemma}[theorem]{Lemma}
\newtheorem{corollary}[theorem]{Corollary}
\newtheorem{remark}[theorem]{Remark}

\newcommand{\E}{\mathbb E}
\newcommand{\Ical}{\mathcal I}
\newcommand{\KL}{D_{\mathrm{KL}}}

\title{A Connectivity Tax Beyond the Scalar Entropy Barrier\\
\large An exact lower certificate for the forest mode constant at $\lambda=1$}
\author{}
\date{2 October 2026}

\begin{document}
\maketitle

\begin{abstract}
Let $\mu(T)$ be the mean size of a uniformly random independent set of a finite tree $T$, and let $\alpha(T)$ be its independence number.  Write
\[
 a_*(1)=\inf_T \frac{\mu(T)}{\alpha(T)},
\]
where the infimum is over finite nonempty trees.  Starting from the entropy--partition decomposition, we prove the strictly stronger universal bound
\[
\boxed{
 a_*(1)\ge
 \frac{\log(34/5)}{6\log 2}
 +\frac{\log(1458/1445)}{48\log 2}
 =0.461191649955056508\ldots>0.461 .
}
\]
The improvement comes from a connectivity tax that is absent from the scalar degree relaxation.  Near scalar equality, almost all excess $A$--degree must be carried by vertices with $d=4$.  If the induced $B$--forest is dense, its edge term already pays a positive cost; if it is sparse, linearly many $d=4$ vertices are isolated in the $B$--forest and can be paired through shared $A$--neighbours.  Under the comparison hard-core measure each such pair is occupied with probability exactly $1/289$, forcing a linear overlap entropy contribution.  The only scalar comparisons needed by the proof are reducible to exact integer arithmetic and are checked by the accompanying script \texttt{verify\_connectivity\_tax\_0461.py}.
\end{abstract}

\section{Statement}

For a finite graph $G$, let
\[
 Z_G=|\Ical(G)|,
 \qquad
 \mu(G)=\frac1{Z_G}\sum_{I\in\Ical(G)}|I|,
 \qquad
 \alpha(G)=\max_{I\in\Ical(G)}|I|.
\]
Set
\[
 c:=\frac{\log(34/5)}{6\log2},
 \qquad
 \gamma:=\frac12-c
 =\frac{\log(20/17)}{6\log2},
\]
and define the new separation constant
\[
 \rho:=\frac{\log(1458/1445)}{12\log2}
 =0.00107676891157466036\ldots .
\]

\begin{theorem}[Connectivity-tax certificate]\label{thm:main}
Every finite nonempty tree $T$ satisfies
\[
 \frac{\mu(T)}{\alpha(T)}
 \ge c+\frac{\rho}{4}.
\]
Consequently
\[
 \boxed{
 a_*(1)\ge c+\frac{\rho}{4}
 =0.461191649955056508\ldots>0.461 .
 }
\]
\end{theorem}

The proof is completely uniform in the order and maximum degree of the tree.

\section{Entropy--partition decomposition}

Fix a maximum independent set $A$ of $T$ and write
\[
 \alpha=|A|,
 \qquad B=V(T)\setminus A,
 \qquad K=T[B],
 \qquad e=|E(K)|.
\]
For $v\in B$ put
\[
 d_v=|N(v)\cap A|.
\]
Since $A$ is maximum, it is maximal, hence $d_v\ge1$ for every $v\in B$.
Because $T$ is a tree,
\begin{equation}\label{eq:edgeidentity}
 \sum_{v\in B}(d_v-1)=\alpha-1-e.
\end{equation}

For $d\ge1$ define
\[
 \ell_d:=\left[\log\frac{5/4}{1+2^{-d}}\right]_+,
 \qquad
 \Delta_d:=\gamma(d-1)-\frac{\ell_d}{2\log2}.
\]
We shall use the following sharpened form of the scalar estimate.

\begin{lemma}[Scalar separation away from $d=4$]\label{lem:scalar}
One has $\Delta_1=\Delta_4=0$, and for every integer $d\ge2$, $d\ne4$,
\[
 \Delta_d\ge \rho(d-1).
\]
\end{lemma}

\begin{proof}
We have
\[
 \ell_3=\log(10/9),
 \qquad
 \ell_4=\log(20/17),
\]
so $\Delta_4=0$ and
\[
 \frac{\Delta_3}{2}
 =\gamma-\frac{\ell_3}{4\log2}
 =\frac{\log(1458/1445)}{12\log2}=\rho.
\]
For $d\ge2$, $d\ne4$, the desired inequality is equivalent to
\begin{equation}\label{eq:scalar_equiv}
 2\ell_d\le(d-1)\ell_3.
\end{equation}
For $d=2$ this is immediate and for $d=3$ it is equality.  For $d=5$ it follows from
\[
 \left(\frac{40}{33}\right)^2
 <\left(\frac{10}{9}\right)^4,
\]
and for every $d\ge6$ from
\[
 2\ell_d<2\log(5/4)
 <5\log(10/9)
 \le(d-1)\ell_3.
\]
The two strict rational comparisons are checked after clearing denominators by the accompanying integer audit.
\end{proof}

We next record exactly the part of the entropy argument that will be used.  Let $J=I\cap B$ for a uniformly random independent set $I$ of $T$.  Put
\[
 w_v=2^{-d_v},
 \qquad
 Z_K(w)=\sum_{J\in\Ical(K)}\prod_{v\in J}w_v.
\]
Let $Q$ be the hard-core measure on $K$ with activities $(w_v)_{v\in B}$.
For $J\in\Ical(K)$ define
\[
 O(J):=\sum_{v\in J}d_v-|N_A(J)|.
\]
Then
\begin{equation}\label{eq:overlap}
 O(J)=\sum_{a\in A}\bigl(k_a(J)-1\bigr)_+,
 \qquad
 k_a(J)=|J\cap N(a)|.
\end{equation}
Counting first over $J$ gives the exact identity
\begin{equation}\label{eq:RA}
 Z_T
 =2^\alpha Z_K(w)\,\E_Q 2^{O(J)},
 \qquad
 R_A:=\log\E_Q2^{O(J)}\ge(\log2)\E_Q O(J).
\end{equation}
The last inequality is Jensen.

For completeness, the entropy identity with the comparison measure of activity $1/4$ is
\[
 2\log2\,\mu(T)
 =\log Z_T-\log Z_K(1/4)+\KL(P\|Q_{1/4}),
\]
where $P$ is the law of $J$ and $Q_{1/4}(J)=4^{-|J|}/Z_K(1/4)$.
Changing the coordinates of $Z_K$ one at a time yields
\[
 \log Z_K(w)-\log Z_K(1/4)
 \ge-\sum_{v\in B}\ell_{d_v}.
\]
Indeed, with all other activities fixed the partition function is $U+xV$ with $0<V\le U$, so lowering a coordinate from $1/4$ to $x\le1/4$ costs at most
$\log[(1+1/4)/(1+x)]$, while increasing a coordinate costs nothing.
Combining these facts and discarding the nonnegative relative entropy gives
\begin{equation}\label{eq:masterlower}
 \mu(T)-c\alpha
 \ge \gamma(1+e)+S+\frac{R_A}{2\log2},
 \qquad
 S:=\sum_{v\in B}\Delta_{d_v}.
\end{equation}
Here we used \eqref{eq:edgeidentity}.

\section{The connectivity tax}

Let
\[
 B_4=\{v\in B:d_v=4\}.
\]
By Lemma~\ref{lem:scalar} and \eqref{eq:edgeidentity},
\begin{equation}\label{eq:n4}
 3|B_4|
 \ge \alpha-1-e-\frac{S}{\rho}.
\end{equation}
Let $B_4^0\subseteq B_4$ be the vertices isolated in $K=T[B]$.
At most $2e$ vertices of $K$ are incident with an edge, hence
\begin{equation}\label{eq:n40}
 |B_4^0|\ge|B_4|-2e.
\end{equation}

For every $a\in A$, let
\[
 q_a=|N(a)\cap B_4^0|.
\]
Pair arbitrarily $2\lfloor q_a/2\rfloor$ of these neighbours.  Let $M$ be the total number of pairs produced over all $a\in A$.  Since
\[
 \sum_{a\in A}q_a=4|B_4^0|,
\]
we have
\begin{align}
 M
 &=\sum_{a\in A}\left\lfloor\frac{q_a}{2}\right\rfloor
 \ge \frac{4|B_4^0|-\alpha}{2} \notag\\
 &\ge
 \frac{\alpha}{6}-\frac23-\frac{14}{3}e
 -\frac{2}{3\rho}S.
 \label{eq:Mlower}
\end{align}
Of course $M\ge0$ as well.

\begin{lemma}[Overlap tax]\label{lem:overlaptax}
With $M$ as above,
\[
 \frac{R_A}{2\log2}\ge\frac{M}{578}.
\]
\end{lemma}

\begin{proof}
Fix $a\in A$.  If $m_a(J)$ of the selected pairs at $a$ have both endpoints in $J$, then
\[
 m_a(J)\le \left\lfloor\frac{k_a(J)}2\right\rfloor
 \le (k_a(J)-1)_+.
\]
Summing over $a$ and using \eqref{eq:overlap},
\[
 O(J)\ge\sum_{\text{selected pairs }\{u,v\}}
 \mathbf 1_{\{u,v\subseteq J\}}.
\]
Every vertex in $B_4^0$ is an isolated vertex of $K$ with activity $2^{-4}=1/16$.
Under $Q$ it is therefore occupied independently with probability $1/17$.
Consequently every selected pair is simultaneously occupied with probability exactly $1/289$, and linearity of expectation gives
\[
 \E_QO(J)\ge\frac{M}{289}.
\]
Now use \eqref{eq:RA}.
\end{proof}

\section{Cancellation and proof of the theorem}

Set
\[
 \theta:=867\rho.
\]
The exact scalar comparisons in the accompanying audit imply
\begin{equation}\label{eq:theta_gamma}
 0<\theta<1,
 \qquad
 \gamma>7\rho.
\end{equation}
Since $M\ge0$ and $M$ also satisfies \eqref{eq:Mlower}, for every $0\le\theta\le1$ we may write
\[
 M\ge \theta\left(
 \frac{\alpha}{6}-\frac23-\frac{14}{3}e
 -\frac{2}{3\rho}S
 \right).
\]
Substituting Lemma~\ref{lem:overlaptax} into \eqref{eq:masterlower} and using this inequality yields
\begin{align*}
 \mu(T)-c\alpha
 &\ge \gamma(1+e)+S
 +\frac{\theta}{578}
 \left(
 \frac{\alpha}{6}-\frac23-\frac{14}{3}e
 -\frac{2}{3\rho}S
 \right)\\
 &=\frac{\rho}{4}\alpha
 +(\gamma-\rho)
 +(\gamma-7\rho)e.
\end{align*}
The $S$--coefficient cancels exactly because
\[
 \frac{867\rho}{578}\frac{2}{3\rho}=1.
\]
By \eqref{eq:theta_gamma}, the two remaining correction terms are positive.  Hence
\[
 \mu(T)-c\alpha\ge\frac{\rho}{4}\alpha,
\]
which proves Theorem~\ref{thm:main}.

\section{Exact scalar audit}

No floating-point comparison is needed for the theorem.  The verifier
\texttt{verify\_connectivity\_tax\_0461.py} checks the following after clearing denominators:
\begin{enumerate}
 \item the $d=5$ comparison
 \[
 (40/33)^2<(10/9)^4;
 \]
 \item the uniform $d\ge6$ comparison
 \[
 (5/4)^2<(10/9)^5;
 \]
 \item $\theta<1$, equivalently
 \[
 (1458/1445)^{289}<16;
 \]
 \item $\gamma>7\rho$, equivalently
 \[
 (20/17)^2>(1458/1445)^7;
 \]
 \item the advertised decimal threshold, equivalently
 \[
 \left(34^8\cdot1458\right)^{1000}
 >2^{22128}\left(5^8\cdot1445\right)^{1000}.
 \]
\end{enumerate}
The last comparison is exactly
\[
 c+\frac{\rho}{4}>\frac{461}{1000}.
\]
The verifier also checks the three rational coefficient identities responsible for the cancellation in the final display.

\section{Current interval and relation to the Bellman attack}

Combining Theorem~\ref{thm:main} with the supplied exact layered-tree witness with branching factors
\[
 (2,1,3,9,6,9,9,9,14,10)
\]
gives the present certified interval
\[
 \boxed{
 0.461191649955056508\ldots
 \le a_*(1)
 <0.473301418185171818.
 }
\]
The upper endpoint is a genuine integer-tree witness.  It should not be confused with the separate fractional Bellman obstruction: the latter concerns a weighted relaxation and is neither a tree witness nor an upper bound on $a_*(1)$.

The significance of the present argument for the Bellman programme is structural.  The old scalar entropy bound becomes sharp only at degree $d=4$.  The proof above shows that a connected tree cannot realize that scalar equality for free: either the induced $B$--forest pays through $e$, or isolated $d=4$ vertices create overlap entropy through shared $A$--neighbours.  Thus the first nonzero ``connectivity correction'' can already be extracted without solving the full Bellman dual.

\end{document}
